\documentclass[11pt]{article}

\usepackage[T1]{fontenc}
\usepackage[utf8]{inputenc}
\usepackage{lmodern}
\usepackage{amsmath,amssymb,amsthm,mathtools}
\usepackage{geometry}
\usepackage{booktabs}
\usepackage{array}
\usepackage{longtable}
\usepackage{hyperref}
\usepackage{enumitem}

\geometry{margin=1in}
\setlist[itemize]{leftmargin=1.5em}
\setlist[enumerate]{leftmargin=1.75em}

\newcommand{\ec}[1]{\texttt{\detokenize{#1}}}
\newcommand{\Fq}{\mathbf{F}_q}
\newcommand{\Z}{\mathbb{Z}}
\newcommand{\Range}[2]{\{#1,\ldots,#2\}}

\newtheorem{observation}{Observation}
\newtheorem{invariant}{Invariant}
\newtheorem{remark}{Remark}

\title{Proof Report for \texorpdfstring{\ec{RefJasminNTT.ec}}{RefJasminNTT.ec}}
\author{Codex Proof Report}
\date{28 April 2026}

\begin{document}
\maketitle

\begin{abstract}
This report documents what was proved in \ec{RefJasminNTT.ec} and how the
proof is structured. The file establishes direct loop-level pRHL refinements
from the extracted Jasmin NTT cores in \ec{Hpoly_extract.M} to the imperative
EasyCrypt reference procedures in \ec{NTT_Fq.NTT}. The forward theorem proves
that \ec{Hpoly_extract.M._poly_ntt} refines \ec{NTT_Fq.NTT.ntt}. The inverse
theorem proves that \ec{Hpoly_extract.M._poly_invntt} refines
\ec{NTT_Fq.NTT.invntt}, with the representation relation strengthened at the
postcondition to account for the final Montgomery-encoded scaling performed by
the Jasmin routine. The proof is not an algebraic proof of the NTT against
\ec{Rq}; it is the concrete extraction-to-imperative-specification layer.
\end{abstract}

\section{Scope of the File}

\ec{RefJasminNTT.ec} is deliberately narrower than the algebraic correctness
development. Its header states the intended boundary: it targets
\ec{extract/Hpoly_extract.ec} and stops at refinement to the imperative
EasyCrypt NTT procedures in \ec{NTT_Fq.ec}. Thus the proof chain covered here is
\[
\text{extracted Jasmin core}
\quad\Longrightarrow\quad
\text{imperative EasyCrypt NTT model}.
\]
The later mathematical step
\[
\text{imperative EasyCrypt NTT model}
\quad\Longrightarrow\quad
\text{abstract ring-theoretic NTT}
\]
is outside the scope of this file.

The two central theorems are:
\begin{align*}
&\ec{poly_ntt_core_ref},\\
&\ec{poly_invntt_core_ref}.
\end{align*}
Two wrapper equivalences are also proved:
\[
\ec{poly_ntt_jazz_ref}
\qquad\text{and}\qquad
\ec{poly_invntt_jazz_ref},
\]
showing that the public extracted wrappers are just inlined calls to the core
procedures. These wrapper proofs are short \ec{proc; inline*; sim} arguments;
the substantial work is in the core refinements.

\section{Target Procedures}

The right-hand programs are the imperative specifications in \ec{NTT_Fq.NTT}.
The forward procedure \ec{NTT.ntt} uses the standard descending layer lengths
\[
128,\ 64,\ 32,\ 16,\ 8,\ 4,\ 2,\ 1
\]
and, for each butterfly, computes
\[
t = \zeta \cdot r_{j+\ell},\qquad
r_{j+\ell}' = r_j - t,\qquad
r_j' = r_j + t.
\]

The inverse procedure \ec{NTT.invntt} uses ascending layer lengths
\[
1,\ 2,\ 4,\ 8,\ 16,\ 32,\ 64,\ 128
\]
and computes
\[
t = r_j,\qquad
r_j' = t + r_{j+\ell},\qquad
r_{j+\ell}' = \zeta \cdot (t - r_{j+\ell}).
\]
After the inverse butterfly layers, \ec{NTT.invntt} performs a final scaling
loop:
\[
r_j' = r_j \cdot \ec{zetas_inv[255]}.
\]

The left-hand programs are the corresponding extracted Jasmin cores:
\[
\ec{Hpoly_extract.M._poly_ntt}
\qquad\text{and}\qquad
\ec{Hpoly_extract.M._poly_invntt}.
\]
These programs manipulate \ec{BArray1024.t} buffers containing 256 signed
32-bit coefficients and call the extracted Montgomery multiplication routine
\ec{Hpoly_extract.M.__fqmul}.

\section{Representation Relations}

\subsection{Base Representation}

The base representation is inherited from \ec{NTT_Fq.ec}. A machine word is
interpreted as a field coefficient by
\[
\ec{word_to_coeff}(w)
  = \operatorname{incoeff}(\ec{W32.to_sint}(w)).
\]
The polynomial readback from a 1024-byte buffer is
\[
\ec{barray256_to_poly}(rp)[i]
  =
\ec{word_to_coeff}(\ec{BArray1024.get32}\; rp\; i).
\]
The relation
\[
\ec{poly_repr}(rp,p)
\]
states that this readback equals the coefficient array \(p\).

The bounded relation used by the final theorems is
\[
\ec{poly_repr_bound}(rp,p,s)
  \equiv
\ec{poly_repr}(rp,p)
\land
\forall i \in \Range{0}{255}.\,
\ec{Fq.bw32}(\ec{BArray1024.get32}\; rp\; i,\ s).
\]
Here \ec{Fq.bw32 w s} records that the signed value of \(w\) is representable
within the signed bound \(2^s\). This is the central device used to rule out
machine-word overflow in additions, subtractions, and products supplied to
\ec{__fqmul}.

\subsection{Pointwise Bounds}

\ec{RefJasminNTT.ec} introduces a more flexible pointwise bounded relation:
\[
\ec{barray256_bound_by}(rp,bd)
  \equiv
\forall i \in \Range{0}{255}.\,
\ec{Fq.bw32}(\ec{BArray1024.get32}\; rp\; i,\ bd(i)).
\]
and
\[
\ec{poly_repr_bound_by}(rp,p,bd)
  \equiv
\ec{poly_repr}(rp,p)\land \ec{barray256_bound_by}(rp,bd).
\]
The pointwise form is necessary because, inside a butterfly loop, some
coefficients in the current stage have already been updated and therefore have
a larger bound than coefficients still waiting to be processed.

The following structural lemmas make this relation usable in loop invariants:
\begin{itemize}
\item \ec{barray256_bound_by_get}: extracts the bound for a concrete index;
\item \ec{barray256_bound_by_weaken}: weakens a pointwise bound;
\item \ec{barray256_bound_by_set32}: preserves a bound after one store;
\item \ec{barray256_bound_by_set32_change}: changes the pointwise bound after
      one store while weakening all untouched entries;
\item \ec{barray256_bound_by_const} and \ec{barray256_bound_by_const_bound}:
      convert between uniform bounds and pointwise constant bounds.
\end{itemize}

\section{Montgomery Table Bridges}

The extracted Jasmin multiplication is Montgomery multiplication. Consequently,
twiddle constants loaded from \ec{jzetas} and \ec{jzetas_inv} are Montgomery
encoded. \ec{RefJasminNTT.ec} imports the concrete table facts from
\ec{NTT_Fq.ec} and packages them as cancellation lemmas.

For the forward table, the key bridge is
\[
\ec{jzetas_get_cancel}:
\quad
\ec{word_to_coeff}(\ec{jzetas}[i])\cdot R^{-1}
  =
\ec{zetas}[i]
\quad (1 \le i < 256).
\]
For the inverse table, the analogous bridge is
\[
\ec{jzetas_inv_get_cancel}:
\quad
\ec{word_to_coeff}(\ec{jzetas_inv}[i])\cdot R^{-1}
  =
\ec{zetas_inv}[i]
\quad (0 \le i < 255).
\]

The final inverse scaling entry is special. The table entry at index \(255\)
is used by the Jasmin final loop. Its cancellation lemma is
\[
\ec{jzetas_inv_255_scale_cancel}:
\quad
\ec{word_to_coeff}(\ec{jzetas_inv}[255])\cdot R^{-1}
  =
\ec{scale255}\cdot R.
\]
This identity explains why the inverse theorem does not end with
\ec{poly_repr_bound res{1} res{2} 16}. The Jasmin final loop produces a
Montgomery-scaled image of the right-hand EasyCrypt result.

\section{Arithmetic Correctness of Extracted Multiplication}

The proof factors every call to \ec{Hpoly_extract.M.__fqmul} through
specialized Hoare and probabilistic Hoare lemmas. At the lowest level,
\ec{fqmul_word_to_coeff_mul_bound_h} proves that if the signed integer product
of the two inputs lies in the range accepted by the signed Montgomery reducer,
then the result has both the correct field value and a 16-bit signed bound:
\[
\begin{aligned}
&\ec{word_to_coeff}(res)
  =
\operatorname{incoeff}(a)\cdot \operatorname{incoeff}(b)\cdot R^{-1},\\
&\ec{Fq.bw32}(res,16).
\end{aligned}
\]
The probabilistic version \ec{fqmul_word_to_coeff_mul_bound_ph} adds the
losslessness side condition, using \ec{Fq.fqmul_ll}.

The product bound used by all butterfly calls is
\[
\ec{fqmul_product_bound_16_24}.
\]
It states that a 16-bit twiddle multiplied by a 24-bit coefficient fits in the
input range required by the Montgomery reduction proof:
\[
-\frac{R}{2}q
\le
\ec{to_sint}(a)\ec{to_sint}(b)
<
\frac{R}{2}q.
\]
The forward and inverse loop invariants are designed precisely so that every
coefficient supplied to \ec{__fqmul} can be weakened to a 24-bit bound, while
every twiddle table entry has a 16-bit bound.

\section{Forward Refinement}

\subsection{Statement}

The forward core theorem is:
\[
\begin{array}{l}
\ec{equiv poly_ntt_core_ref :}\\
\quad
\ec{Hpoly_extract.M._poly_ntt ~ NTT_Fq.NTT.ntt :}\\
\quad
\ec{NTT_Fq.poly_repr_bound rp\{1\} r\{2\} 16 /\\}\\
\quad
\ec{zetas\{2\} = NTT_Fq.zetas}\\
\quad
\ec{==> NTT_Fq.poly_repr_bound res\{1\} res\{2\} 24.}
\end{array}
\]

The precondition says that the Jasmin input buffer represents the right-hand
coefficient array and all input words are initially 16-bit bounded. The
postcondition says that the two procedures return represented arrays, with a
24-bit output bound. The 24-bit bound is not accidental: each forward stage can
increase the coefficient bound by one bit, and there are eight butterfly
layers.

\subsection{Stage Bookkeeping}

Forward loop control is summarized by:
\[
\ec{fwd_len_ok}(len)
\]
for the allowed values \(128,64,\ldots,1,0\), and
\[
\ec{fwd_stage_bound}(len)
\]
for the uniform stage bound:
\[
\begin{array}{c|ccccccccc}
len & 128&64&32&16&8&4&2&1&0\\
\midrule
\ec{fwd_stage_bound}(len)&16&17&18&19&20&21&22&23&24.
\end{array}
\]
The table counter is tracked by
\[
\ec{fwd_zbase}(len),
\]
whose values correspond to the first twiddle index of the current layer:
\[
0,1,3,7,15,31,63,127,255.
\]
The lemmas \ec{fwd_stage_next}, \ec{fwd_len_next_ok}, and
\ec{fwd_zbase_next} prove that these values evolve exactly as the two programs
update \ec{len} and \ec{zetasctr}.

\subsection{Nested Loop Invariants}

The proof uses three nested pRHL \ec{while} invariants matching the three nested
loops of the programs.

\paragraph{Outer invariant.}
At the start of each layer:
\[
\begin{aligned}
&\ec{poly_repr}(rp_1,r_2),\\
&\ec{barray256_bound_by}(rp_1,\lambda i.\ec{fwd_stage_bound}(len_1)),\\
&\ec{zetasp_1 = jzetas},\qquad
\ec{zetas_2 = NTT_Fq.zetas},\\
&\ec{zetasctr_1 = zetasctr_2},\qquad
\ec{len_1 = len_2},\\
&\ec{fwd_len_ok}(len_1),\qquad
\ec{zetasctr_1 = fwd_zbase}(len_1).
\end{aligned}
\]
This invariant aligns the two programs at layer boundaries and asserts the
uniform bound for the current stage.

\paragraph{Middle invariant.}
Inside a layer, \ec{start} marks how much of the array has already been
processed. The bound is
\[
\ec{fwd_middle_bound}(sz,start,i)
  =
\begin{cases}
sz+1,&0\le i<start,\\
sz,&\text{otherwise}.
\end{cases}
\]
Thus coefficients in completed blocks have grown by one bit, while the rest of
the layer remains at the original stage bound.

\paragraph{Inner invariant.}
Inside one block, the inner-loop bound is
\[
\ec{fwd_inner_bound}(sz,\ell,start,j,i).
\]
It gives bound \(sz+1\) to coefficients in completed butterflies and bound
\(sz\) to the two current butterfly inputs. The lemmas
\ec{fwd_inner_current}, \ec{fwd_inner_weaken_unchanged}, and
\ec{fwd_inner_exit_to_middle} justify, respectively, reading the current input
bounds, preserving untouched entries after a store, and converting the inner
bound back to the middle bound when the block ends.

\subsection{Butterfly Step}

The forward butterfly proof is split into two conceptual pieces.

First, \ec{forward_twiddle_product_bound_call_ph} verifies the call to
\ec{__fqmul}. It uses:
\begin{itemize}
\item \ec{jzetas_get_cancel} to convert the Montgomery table entry into the
      abstract \(\zeta\);
\item \ec{forward_twiddle_product_bound_from_inner} to prove the machine-word
      product bound required by \ec{__fqmul};
\item \ec{fqmul_word_to_coeff_mul_bound_ph} to obtain both semantic
      correctness and a 16-bit result bound.
\end{itemize}

Second, \ec{forward_butterfly_step_bound_by} verifies the two stores:
\[
r_{j+\ell}\leftarrow r_j-t,\qquad
r_j\leftarrow r_j+t.
\]
The field-level equality is obtained through \ec{word_to_coeff_add} and
\ec{word_to_coeff_sub}; their preconditions are exactly the pointwise
\ec{bw32} bounds carried in the invariant. The output bound is raised from
\(sz\) to \(sz+1\) on the two updated positions using
\ec{forward_butterfly_bounds}.

\subsection{Completion of the Forward Proof}

At inner-loop exit, the proof derives \(j=start+\ell\) and converts
\ec{fwd_inner_bound} to \ec{fwd_middle_bound}. At middle-loop exit, it derives
\(start=256\), upgrades the whole layer to \(sz+1\), and uses
\ec{fwd_stage_next} to re-enter the outer invariant at the next value
\(\ell/2\). At outer-loop exit, \(\ell=0\), and
\ec{fwd_stage_bound}(0) is 24, giving the final postcondition
\[
\ec{poly_repr_bound}(res_1,res_2,24).
\]

\section{Inverse Refinement}

\subsection{Statement}

The inverse core theorem is:
\[
\begin{array}{l}
\ec{equiv poly_invntt_core_ref :}\\
\quad
\ec{Hpoly_extract.M._poly_invntt ~ NTT_Fq.NTT.invntt :}\\
\quad
\ec{NTT_Fq.poly_repr_bound rp\{1\} r\{2\} 16 /\\}\\
\quad
\ec{zetas_inv\{2\} = NTT_Fq.zetas_inv}\\
\quad
\ec{==> NTT_Fq.poly_repr_bound res\{1\}}\\
\quad\qquad
\ec{(NTT_Fq.array256_mont res\{2\}) 16.}
\end{array}
\]

The postcondition is intentionally strengthened relative to a plain
coefficient representation. The right-hand program multiplies by
\ec{zetas_inv[255]} at the abstract coefficient level, whereas the left-hand
Jasmin code calls Montgomery multiplication with the concrete table entry
\ec{jzetas_inv[255]}. The table bridge proves that this machine operation
produces the Montgomery-scaled coefficient
\[
r_j\cdot \ec{scale255}\cdot R.
\]
Therefore the final represented polynomial is
\[
\ec{array256_mont}(res_2),
\]
not \(res_2\) itself.

\subsection{Inverse Stage Bookkeeping}

The inverse layer control is summarized by:
\[
\ec{inv_len_ok}(len)
\]
for the values \(1,2,\ldots,128,256\), and
\[
\ec{inv_stage_bound}(len)
\]
with values:
\[
\begin{array}{c|ccccccccc}
len&1&2&4&8&16&32&64&128&256\\
\midrule
\ec{inv_stage_bound}(len)&16&17&18&19&20&21&22&23&24.
\end{array}
\]
The inverse table counter is described by
\[
\ec{inv_zbase}(len)
 =
0,128,192,224,240,248,252,254,255.
\]
The lemmas \ec{inv_stage_next}, \ec{inv_len_next_ok}, and
\ec{inv_stage_zbase_exit} prove that these values match the update
\(\ell\leftarrow 2\ell\) and the inverse table traversal.

\subsection{Inverse Butterfly Layers}

The inverse proof reuses the same three-level invariant shape as the forward
proof, replacing \ec{fwd_*} operators with \ec{inv_*}. The important semantic
difference is the butterfly:
\[
r_j' = r_j+r_{j+\ell},\qquad
r_{j+\ell}' = \zeta\cdot(r_j-r_{j+\ell}).
\]
The extracted code computes the difference in a temporary word and then calls
\ec{__fqmul}. The proof therefore needs two additional facts:
\begin{itemize}
\item \ec{inverse_twiddle_product_bound_from_inner}: the word difference
      \(rp[j]-rp[j+\ell]\) is bounded by \ec{inv_stage_bound len + 1};
\item \ec{inverse_twiddle_product_call_bound_from_inner}: after weakening this
      difference to a 24-bit bound and combining it with a 16-bit inverse
      twiddle, the product satisfies the Montgomery reducer input bound.
\end{itemize}

The call to \ec{__fqmul} is discharged by
\ec{inverse_twiddle_product_bound_call_ph}. The post-call store reasoning is
handled by \ec{inverse_butterfly_step_bound_by}. The array update shape of the
right-hand specification differs slightly from the extracted code order, so
\ec{inverse_spec_updateE} proves the extensional equality between the generated
store expression and the compact mathematical update:
\[
(p[j\leftarrow p[j]+p[j+\ell]])[j+\ell\leftarrow
\zeta(p[j]-p[j+\ell])].
\]

\subsection{Final Scaling Loop}

After the eight inverse butterfly layers, the coefficients have a uniform
24-bit bound. The final loop applies \ec{__fqmul} at each index with
\ec{jzetas_inv[255]}. This loop has its own representation and bound
invariant:
\[
\ec{poly_repr}(rp_1,\ec{final_mont_array}(r_2,j_2))
\land
\ec{barray256_bound_by}(rp_1,\ec{final_bound}(j_1)).
\]
The ghost array
\[
\ec{final_mont_array}(p,j)[i]
 =
\begin{cases}
p[i]\cdot R,&0\le i<j,\\
p[i],&\text{otherwise}
\end{cases}
\]
records that the prefix already processed by the final loop has been
Montgomery-scaled. The corresponding bound function is
\[
\ec{final_bound}(j,i)
 =
\begin{cases}
16,&0\le i<j,\\
24,&\text{otherwise}.
\end{cases}
\]
Thus processed entries are known to be fresh 16-bit Montgomery multiplication
outputs, while unprocessed entries still carry the 24-bit bound inherited from
the inverse butterfly layers.

The key final-loop lemmas are:
\begin{itemize}
\item \ec{final_mont_array_start}: before the loop, the ghost array is the
      unscaled right-hand result;
\item \ec{final_mont_array_setE}: one iteration advances the Montgomery-scaled
      prefix by one entry;
\item \ec{final_mont_array_exit}: after \(j=256\), the ghost array equals
      \ec{NTT_Fq.array256_mont p};
\item \ec{inverse_final_product_bound_call_ph}: the \ec{__fqmul} call at
      index \(255\) computes \(coeff\cdot\ec{scale255}\cdot R\) and returns a
      16-bit word;
\item \ec{inverse_final_product_bound_from_inv}: every final-loop product
      satisfies the signed Montgomery reducer range because the scale entry is
      16-bit and the current unprocessed coefficient is at most 24-bit.
\end{itemize}

\subsection{Completion of the Inverse Proof}

The inverse theorem is organized as two pRHL loops after \ec{proc; wp}. First,
the proof discharges the final scaling loop because it appears syntactically
last and is therefore encountered first by weakest precondition reasoning.
Then it proves the nested inverse butterfly loops.

At the exit of the inverse butterfly loops, the invariant gives
\[
\ec{poly_repr}(rp_1,r_2)
\land
\ec{barray256_bound_by}(rp_1,\lambda i.24).
\]
This state is converted into the precondition of the final-loop invariant by
\ec{final_mont_array_start} and \ec{final_bound_start}. At the exit of the
final loop, \(j=256\); \ec{final_mont_array_exit} converts the representation
target to \ec{NTT_Fq.array256_mont res{2}}, and \ec{final_bound} gives the
uniform 16-bit postcondition.

\section{Machine-Word Overflow Argument}

The overflow argument is not a separate after-the-fact side condition. It is
threaded through all pRHL invariants.

The core facts are:
\begin{enumerate}
\item All input coefficients begin with a 16-bit bound.
\item Every butterfly layer can increase updated coefficient bounds by at most
      one bit, as proved by \ec{forward_butterfly_bounds} and
      \ec{inverse_butterfly_bounds}.
\item There are eight butterfly layers, so a coefficient entering a
      Montgomery multiplication is never worse than 24-bit.
\item All twiddle-table entries used in multiplication are 16-bit:
      \ec{jzetas_bound16} and \ec{jzetas_inv_bound16}.
\item A 16-bit twiddle times a 24-bit coefficient satisfies the reducer
      precondition by \ec{fqmul_product_bound_16_24}.
\item Every result of \ec{__fqmul} is reduced back to a 16-bit word.
\end{enumerate}

This is why the forward theorem ends with a 24-bit bound, while the inverse
theorem ends with a 16-bit bound on the Montgomery-scaled output: the inverse
has the same eight-layer growth to 24 bits, followed by a final multiplication
that reduces every entry back to 16 bits.

\section{Proof Engineering Pattern}

The proof uses a consistent engineering pattern throughout:
\begin{enumerate}
\item Define a small stage-indexing model for the loop counter values.
\item Define a pointwise bound function that distinguishes completed,
      current, and untouched coefficients.
\item Prove arithmetic transition lemmas for that bound function.
\item Verify the extracted Montgomery multiplication call with a dedicated
      \ec{phoare} lemma.
\item Verify the subsequent stores by combining field representation lemmas
      with machine-word add/subtract bounds.
\item Close each loop by converting the inner invariant to the middle
      invariant, then the middle invariant to the next outer stage invariant.
\end{enumerate}

This structure keeps the pRHL proof direct: the left and right programs are
related at the actual loop level, without replacing the extracted code by a
separate hand-written model.

\section{Checked Artifacts}

The relevant checked artifacts are:
\begin{itemize}
\item \ec{RefJasminNTT.ec}, containing the pRHL refinements and the supporting
      representation and bound lemmas;
\item \ec{NTT_Fq.ec}, containing the imperative NTT procedures, table
      definitions, table bridges, and base representation relation;
\item \ec{extract/Hpoly_extract.ec}, containing the extracted Jasmin
      procedures and concrete twiddle tables.
\end{itemize}

The proof was checked with:
\begin{verbatim}
easycrypt compile RefJasminNTT.ec -I .
easycrypt compile NTT_Fq.ec -I .
\end{verbatim}
and the proof-marker scan
\begin{verbatim}
rg -n "admit|abort|axiom" RefJasminNTT.ec NTT_Fq.ec
\end{verbatim}
reported no proof holes in \ec{RefJasminNTT.ec}. The only reported matches were
the word ``axiom'' in an explanatory comment in \ec{NTT_Fq.ec}.

\section{Summary}

\ec{RefJasminNTT.ec} proves the direct loop-level extraction refinements for
both NTT cores. The forward proof establishes:
\[
\ec{Hpoly_extract.M._poly_ntt}
\sim
\ec{NTT_Fq.NTT.ntt}
\]
under a 16-bit input representation relation, producing a 24-bit represented
output. The inverse proof establishes:
\[
\ec{Hpoly_extract.M._poly_invntt}
\sim
\ec{NTT_Fq.NTT.invntt}
\]
under the analogous inverse-table precondition, producing a 16-bit represented
output for \ec{array256_mont} of the right-hand result.

The main technical content is the pointwise machine-word bound infrastructure.
It rules out overflow locally at every butterfly and at every Montgomery
multiplication call, while simultaneously preserving the field-level
representation relation needed for the pRHL simulation.

\end{document}
